\section{Vector Fields}

In this section, we introduce vector fields on surfaces and use them to develop an alternative proof of the Gauss--Bonnet theorem. This proof relies on another remarkable local-to-global result, the Poincar'e--Hopf Index Theorem, which relates the local behavior of a vector field near its zeros to the topology of the underlying surface.



\subsection{Poincar\'e Index Theorem}

A \emph{tangent vector field} on a surface is an assignment of a tangent 
vector to each point of the surface. A vector field is said to be \emph{smooth} 
if this assignment varies smoothly from point to point. 
See \figref{vector-field} for an example.

Points where the assigned vector is the zero vector are special.
 We assume, these zeros occur as isolated points. 
 Remarkably, the arrangement of these zeros encodes 
 global topological information about the surface.



 \begin{figure}[htb]
         \centering
         \includegraphics[width=5cm]{background/vector-field}
	\caption{An example of a vector field.
	\label{fig:vector-field}}
 \end{figure}
 
 Let $X$ be a smooth vector field and let $z$ be an isolated zero of $X$. Choose a sufficiently small simple closed curve $\gamma$ surrounding $z$ and containing no other zeros of $X$. Since $X$ is nonzero on $\gamma$, we may normalize the vectors to obtain a map

$$\widehat{X} : \gamma \to \Sp^1,
\qquad
\widehat{X}(p)=\frac{X(p)}{|X(p)|}.
$$

As one traverses $\gamma$ once in the counterclockwise direction, 
the normalized vectors may rotate around the unit circle several times. 
The number of times that $\widehat{X}$ winds around $\mathbb{S}^1$ is 
 the \emph{index} of $X$ at $z$.
 
 \begin{definition}
Let $z$ be an isolated zero of a smooth vector field $X$. The \emph{index} of $X$ at $z$, 
denoted $i(X,z)$, is the winding number of the map
$$\widehat{X} : \gamma \to \Sp^1.$$

Equivalently, $i(X,z)$ measures the net number of rotations made by the vectors of $X$ as one traverses a small loop enclosing $z$.
\end{definition}
 
The definition of index may seem mysterious at first, but it admits a useful interpretation
 in terms of the Euler characteristic. Recall the proof of 
 Euler's formula in \thmref{euler}. 
 Recall the proof of Euler's formula in \thmref{euler}, where we assigned a positive charge to each vertex, a negative charge to each edge, and a positive charge to each face of a triangulation. The total charge was
$$|V|-|E|+|F|=\chi(S).$$

Now imagine allowing these charges to flow along the trajectories determined by a vector field $X$.
 As the charges move, they accumulate at the zeros of the vector field. 
 The index of a zero can be viewed as the amount of net charge that collects there. 
 From this perspective, the Poincar'e--Hopf theorem states that the total charge accumulated at all zeros is equal to the Euler characteristic of the surface.
 

 
 
\begin{theorem}[Poincar\'e Index Theorem] \label{thm:poincare-index}

If $S$ is a regular triangulated surface and $X$ a vector field on $S$ with isolated zeros,
the Euler characteristic of $S$ is 
$$\chi(S)=\sum_{z\in \text{zeros}}i(X,z).$$
\end{theorem}

\begin{proof}

We begin by subdividing and perturbing the triangulation, if necessary, 
so that each zero of $X$ lies in the interior of a unique face
We also assume that $X$ is transverse to every edge of the triangulation. 
Establishing that these assumptions can always be achieved requires additional work,
 which we omit.


As in the proof of Euler's formula, assign a charge of $+1$ to each vertex, 
a charge of $-1$ to each edge, and a charge of $+1$ to each face. 
We then allow these charges to flow along the trajectories of the vector field $X$.


Consider a face that does not contain a zero of $X$.
 Since the vector field is nonvanishing on this face, the charges 
 entering and leaving the face balance, and the net contribution 
 of the face to the total charge is zero. Consequently, any nonzero 
 contribution to the Euler characteristic must arise from faces containing 
 zeros of the vector field.
 
 For a face containing an isolated zero $z$, the net charge that accumulates at $z$ is precisely the index $i(X,z)$. Summing over all such faces, we find that the total charge is
$$\sum_z i(X,z).$$
Because the flow of charge preserves the total amount of charge, this sum must equal the initial total charge of the triangulation, namely
$$\chi(S).$$
Therefore,
$$\chi(S)=\sum_z i(X,z),$$
completing the proof.
\end{proof}


\subsection{Banchoff's Proof}
Another interpretation is presented by Banchoff in \cite{banchoff_critical_1970}.
Given a surface $S$ embedded in $\R^3$ consider a line $\ell$ and define a 
linear height function $h:M \to \R$
by projecting points of $\R^3$ onto  $\ell$. 

A point $p$ on $S$
is a \emph{critical point} of $h$ if the tangent plane at $p$ is perpendicular to $\ell$.
All points on $M$ that are not critical are \emph{ordinary point}.
Critical points can be classified into three categories, maxima, minima, and saddle points.

The height function induces a vector field on $S$ by flowing in the direction
of greatest decreasing height.
Moreover, $i(p,h)=1$ if $p$ is a  local maximum or minimum and $i(p,h)=-1$ is $p$ is a saddle point.
An example is shown in \figref{torus-total}.


For an ordinary point $q$, the tangent plane at $q$ is not horizontal. 
Now consider the plane perpendicular to $\ell$ and through $q$, this plane
divides a `small' neighborhood $U$ of $q$ into two pieces and intersects
a `small' circle about $q$ in two points.
this observation leads to an alternative description of the index.
Let $\mathcal{C}$ denote the number of points a plane through $p$ perpendicular
to $h$ meets a `small' circle around $p$, then 
$$i(p,h)=1-\frac{1}{2}\mathcal{C}.$$

For a  polyhedral surface, every point in the interior of a face or edge
is ordinary, since we assumed that no face or edge is perpendicular to our line $\ell.$

Let $v$ be a vertex of a triangulated surface.
The union of all triangles incident to $v$ is called the \emph{star} of $v$,
 denoted $\mathrm{star}(v).$ A triangle in $\mathrm{star}(v)$
 is a \emph{straddle}
 vertex can be any of the three types of critical points.
So, for each vertex $v$, we cut out the triangles incident to it. This collection of triangles
is called the \emph{star}$(v)$ of $v$. 
A triangle in $\mathrm{star}(v)$ is said to \emph{straddle} the plane through 
$v$ perpendicular to $\ell$ if one of its vertices lies on one side of the plane 
and another lies on the opposite side. Equivalently, among the three vertices 
of the triangle, the height of $v$ is neither the largest nor the smallest. 
In this case, $v$ is called the \emph{middle vertex} of the triangle with respect to $\ell$.

The index of $v$ can be computed combinatorially as
$$i(v,\ell)=1-\frac{1}{2} \left(\text{number of triangles with } v \text{ middle for } \ell\right).$$



T compute the Euler characteristic, note that for a triangulated surface, 
each face is a triangle and has three edges and each
edge is incident to two faces, so
\begin{equation} \label{eqn:faces-edges}
	3F=2E.
\end{equation}

Summing the indices over all vertices yields
\begin{align}
\sum_{v\in M}i(v,\ell) &=\sum_{v\in M} \left(1-\frac{1}{2}(\text{number of triangles with } v \text{ middle for } \ell)\right),  \\ \nonumber
       &=V - \frac{1}{2}\sum_{v\in M} \left(\text{number of triangles with } v \text{ middle for } \ell\right), \\ \nonumber
       &=V-\frac{1}{2} F \text{ ,because each triangle has exactly one middle vertex for  } \ell,\\ \nonumber
       &= V-\frac{1}{2}(2E-2F), \text{by \eqnref{faces-edges}}\\ \nonumber
       &=V-E+F. \nonumber 
\end{align}



 \begin{figure}[htb]
         \centering
        \begin{subfigure}[b]{0.3\textwidth}
         \includegraphics[width=\textwidth]{chapter-3/torus}
         \subcaption{The tours a line $\ell$ with four critical points.}
 	 \label{fig:torus-ell}
       \end{subfigure}\\
         \begin{subfigure}[b]{0.21\textwidth}
         \includegraphics[width=\textwidth]{chapter-3/torus-top}
         \subcaption{Local to $p_1$.}
          \label{fig:max}
         \end{subfigure}
          \begin{subfigure}[b]{0.21\textwidth}
         \includegraphics[width=\textwidth]{chapter-3/saddle1}
         \subcaption{Local to $p_2$.}
          \label{fig:saddle1}
         \end{subfigure}
         \begin{subfigure}[b]{0.21\textwidth}
         \includegraphics[width=\textwidth]{chapter-3/saddle2}
         \subcaption{Local to $p_3$.}
          \label{fig:saddle2}
         \end{subfigure}
         \begin{subfigure}[b]{0.21\textwidth}
         \includegraphics[width=\textwidth]{chapter-3/torus-bottom}
         \subcaption{Local to $p_4$.}\label{fig:min}
         \end{subfigure}
		\caption{(\subref{fig:torus-ell}) A tours and line $\ell$. There are four critical points.
 		(\subref{fig:max}) The critical point $p_1$ is a maximum and $i(p_1,h)=1$.
		(\subref{fig:saddle1}) The critical point $p_2$ is a saddle and $i(p_1,h)=-1$.
		(\subref{fig:saddle2}) The critical point $p_3$ is also a saddle and $i(p_1,h)=-1$.
		(\subref{fig:min}) The critical point $p_4$ is a minimum and $i(p_1,h)=1$.
 		\label{fig:torus-total}}
 \end{figure}

With the Poincar'e--Hopf Index Theorem in hand, we are now ready to prove the Gauss--Bonnet theorem for closed surfaces (\thmref{g-b-nb}).


\begin{proof}

We begin by relating curvature to the index of a height function. Let $U$ be a convex open subset of a smooth surface, and let $g\to \mathbb{S}^2$ denote the Gauss map. Let $dA$ denote the area element on the unit sphere. Recall that
$$\int_{\mathbb{S}^2} dA = 4\pi.$$

The curvature of $U$ can be interpreted as the area of its image under the Gauss map. 
If $g$ is one-to-one on $U$, then

$$\int_{\mathbb{S}^2} F_{g(U)}(\ell),d\omega,$$


where $F_{g(U)}$ is the characteristic function of $g(U)$, defined by


$$F_{g(U)}(\ell)=
\begin{cases}
1, & \ell\in g(U),\\
0, & \ell\notin g(U).
\end{cases}$$


A point $\ell\in \mathbb{S}^2$ determines a height function on $U$. The height function has a critical point at $p\in U$ precisely when $\ell$ is parallel to the unit normal at $p$, that is, when $\ell=g(p)$ or $\ell=-g(p)$. Furthermore, if $U$ is sufficiently small so that $g(U)$ contains no pair of antipodal points, then
the curvature can then be expressed as 
$$K(U)=\int_{\Sp^2}F_{g(U)}(\ell)d\omega =\frac{1}{2}\int_{\Sp^2}\sum_{p\in U} i(p,\ell)d\omega.$$
The right-hand term is well defined for orientation reversing sets $W$, only a negative is needed.

For a neighborhood $U$ where $g$ is not one-to-one, we decompose $U$ into a countable
disjoint union of sets $U_i$ where $g$ is one-to-one with a set $V=U-\cup_{i=1}^\infty U_i$
such that $g(V)$ has measure zero on $\Sp^2$ and therefore
$$K(U)=\int_{U} KA$$ is the sum converges, this agrees with 
$$K(U)=\frac{1}{2}\int_{\Sp^2}\sum_{p\in U} i(p,\xi)d\omega,$$
where the integrand is well defined almost everywhere since almost every $\xi$
has finitely many critical points.

For a smooth surface, we define the point function $K(p)$ as the limit as the diameter of 
and open set $U$ with $p\in U$ goes to zero. Thus, we have the following.

\begin{align}
K(S) &=\frac{1}{2}\int_{\Sp^2}\sum_{p\in U} i(p,\xi)d\omega,  \\ \nonumber
       &=\frac{1}{2}\int_{\Sp^2}\chi(S) \text{ , by the Poincar\'e Index theorem}\\ \nonumber
       &=\frac{1}{2}\chi(S)\int_{\Sp^2}d\omega, \\ \nonumber
       &=2\pi\chi(S). \nonumber 
\end{align}
This proves the Gauss--Bonnet theorem.

\end{proof}




\noindent \textbf{Exercises}


\begin{enumerate}
	\item 
	
	\item 
	
	\item 
	
\end{enumerate}

\pagebreak
